% Provenance (never printed): see sources/outline_sources.md and sources/book1/BRIEFS.md (g7/01).
\chapter{Functions}\label{ch:g7:functions}

Lina's family is repainting the house, and her \omterm{def:g4:programs-and-operating-system:program}{program} computes the area of
each of the five rooms: the same three lines, copied five times with other
numbers. Then she notices a mistake in those three lines. Now she must correct
it five times, and hope she has not missed one. If the three lines had a name,
she could write them once and use them five times, and a mistake would be fixed
in one place.

\begin{recall}
A \omterm{def:g6:variables-and-input:value}{variable} is a name for a \omterm{def:g6:variables-and-input:value}{value} (\cref{ch:g6:variables-and-input}); \code{print}
shows a \omterm{def:g6:variables-and-input:value}{value} on the screen (\cref{ch:g6:first-python-program}); \code{if}
chooses between blocks (\cref{ch:g6:if-else}).
\end{recall}

\section{Defining and calling a function}

\begin{definition}[Function and function call]\label{def:g7:functions:function}
A \emph{function}\index{function} is a block of \omterm{def:g1:instructions-in-order:instruction}{instructions} with a name,
written once and run whenever it is needed. It is \emph{defined} with
\code{def}, the name, brackets and a colon, followed by an indented block, its
\emph{body}. A \emph{function call}\index{function call} is the name followed by
brackets: it runs the body.
\end{definition}

\Needspace*{12\baselineskip}
\omcode{code/grade-7/01-functions/listings/greet.py}{1}{9}{A \omterm{def:g7:functions:function}{function} defined
once and called three times.}\label{lst:g7:functions:greet}

\Needspace*{11\baselineskip}
\omoutput{code/grade-7/01-functions/outputs/greet.txt}{The body ran three times,
once for each call.}\label{out:g7:functions:greet}

Lines 1 to 3 only \emph{define} the \omterm{def:g7:functions:function}{function}: Python remembers it and prints
nothing. That is why \code{Before the calls} comes first. A \omterm{def:g7:functions:function}{function} runs only
when it is called. Two empty lines after a \omterm{def:g7:functions:function}{function} are the custom in Python:
they make each \omterm{def:g7:functions:function}{function} easy to see.

\section{Parameters and arguments}

Most \omterm{def:g7:functions:function}{functions} need \omterm{def:g6:variables-and-input:value}{values} to work on. The area of a room needs its width and
its length.

\begin{definition}[Parameter and argument]\label{def:g7:functions:parameter}
A \emph{parameter}\index{parameter} is a \omterm{def:g6:variables-and-input:value}{variable} named between the brackets
of a \omterm{def:g7:functions:function}{function}'s definition. An \emph{argument}\index{argument} is the \omterm{def:g6:variables-and-input:value}{value}
given for it in a call: in \code{area(3, 4)}, the parameter \code{width} takes
the argument 3, and \code{length} takes 4.
\end{definition}

\begin{definition}[Return value]\label{def:g7:functions:return-value}
The \omterm{def:g1:instructions-in-order:instruction}{instruction} \code{return} ends the call and gives a \omterm{def:g6:variables-and-input:value}{value} back to the place
where the \omterm{def:g7:functions:function}{function} was called: the \emph{return value}\index{return value}. The
call can then be used like that \omterm{def:g6:variables-and-input:value}{value}, in a calculation or an \omterm{def:g6:variables-and-input:assignment}{assignment}.
\end{definition}

\Needspace*{11\baselineskip}
\omcode{code/grade-7/01-functions/listings/area.py}{1}{8}{A \omterm{def:g7:functions:function}{function} with two
\omterm{def:g7:functions:parameter}{parameters} and a \omterm{def:g7:functions:return-value}{return value}.}\label{lst:g7:functions:area}

\Needspace*{7\baselineskip}
\omoutput{code/grade-7/01-functions/outputs/area.txt}{Each call is replaced by
the \omterm{def:g6:variables-and-input:value}{value} it returns.}\label{out:g7:functions:area}

\begin{omfigure}
\begin{tikzpicture}[x=1cm, y=1cm, font=\footnotesize]
  \node[draw=black!70, fill=ocYellow!35, rounded corners=6pt, minimum width=3.4cm, minimum height=1.6cm, align=center] (m) at (0,0) {\code{area}\\[2pt]\scriptsize\code{return width * length}};
  \node[draw=black!50, fill=ocBlue!12, minimum width=1.2cm] (w) at (-4.0,0.45) {\code{width} = 3};
  \node[draw=black!50, fill=ocBlue!12, minimum width=1.2cm] (l) at (-4.0,-0.45) {\code{length} = 4};
  \node[draw=black!50, fill=ocMint!30, minimum width=1.0cm] (r) at (3.7,0) {12};
  \draw[-{Stealth[length=2mm]}, thick] (w.east) -- (m.west |- w.east);
  \draw[-{Stealth[length=2mm]}, thick] (l.east) -- (m.west |- l.east);
  \draw[-{Stealth[length=2mm]}, thick] (m.east) -- (r.west);
  \node[above, black!60] at (-4.0,0.9) {parameters, given the arguments};
  \node[above, black!60] at (3.7,0.4) {return value};
\end{tikzpicture}

{\small A \omterm{def:g7:functions:function}{function} is like a machine: the \omterm{def:g7:functions:parameter}{arguments} go in through the
\omterm{def:g7:functions:parameter}{parameters}, and the \omterm{def:g7:functions:return-value}{return value} comes out.}
\end{omfigure}

\omterm{def:g7:functions:parameter}{Arguments} are given to \omterm{def:g7:functions:parameter}{parameters} in order: the first \omterm{def:g7:functions:parameter}{argument} to the first
\omterm{def:g7:functions:parameter}{parameter}, and so on. Swapping them is not an error for Python, but the
result is wrong:

\Needspace*{9\baselineskip}
\omcode{code/grade-7/01-functions/listings/two_params.py}{1}{6}{The order of the
\omterm{def:g7:functions:parameter}{arguments}.}\label{lst:g7:functions:two-params}

\Needspace*{6\baselineskip}
\omoutput{code/grade-7/01-functions/outputs/two_params.txt}{The second call gave
11 to \code{name} and \code{"Noor"} to \code{age}.}\label{out:g7:functions:two-params}

A \omterm{def:g7:functions:function}{function} can call another \omterm{def:g7:functions:function}{function}. Here, \code{largest3} finds the largest
of three numbers with two calls to \code{largest2}. A \code{return} inside an
\code{if} ends the call at once, so line 4 runs only when \code{a > b} is false.

\omcode{code/grade-7/01-functions/listings/largest.py}{1}{13}{A \omterm{def:g7:functions:function}{function} built
from another.}\label{lst:g7:functions:largest}

\Needspace*{7\baselineskip}
\omoutput{code/grade-7/01-functions/outputs/largest.txt}{The largest of 4 and 9;
of 4, 9 and 2; of 7, 3 and 5.}\label{out:g7:functions:largest}

\begin{omfigure}
\begin{tikzpicture}[x=1cm, y=1cm, font=\footnotesize]
  \node[draw=black!60, fill=black!5, rounded corners=3pt, anchor=west, minimum width=4.4cm] (c) at (0,0) {\code{print(largest3(4, 9, 2))}};
  \node[draw=black!60, fill=ocYellow!25, rounded corners=3pt, anchor=west, minimum width=4.4cm] (l3) at (0,-1.9) {\code{largest3}: $a = 4$, $b = 9$, $c = 2$};
  \node[draw=black!60, fill=ocBlue!10, rounded corners=3pt, anchor=west] (x1) at (6.6,-1.2) {\code{largest2(4, 9)} returns 9};
  \node[draw=black!60, fill=ocBlue!10, rounded corners=3pt, anchor=west] (x2) at (6.6,-2.6) {\code{largest2(9, 2)} returns 9};
  \draw[-{Stealth[length=2mm]}, thick] (0.8,-0.3) -- (0.8,-1.6) node[midway, left] {1. call};
  \draw[-{Stealth[length=2mm]}, thick, ocRed] (2.6,-1.6) -- (2.6,-0.3) node[midway, right] {4. returns 9};
  \draw[-{Stealth[length=2mm]}, thick, ocBlue!70!black] (l3.east) -- node[above, sloped] {2.} (x1.west);
  \draw[-{Stealth[length=2mm]}, thick, ocBlue!70!black] (l3.east) -- node[below, sloped] {3.} (x2.west);
\end{tikzpicture}

{\small The journey of a call. The caller waits while \code{largest3} runs;
\code{largest3} itself waits for each call to \code{largest2}. Each \omterm{def:g6:variables-and-input:value}{value} comes
back to the place where its call was written.}
\end{omfigure}

\section{\texorpdfstring{\texttt{return}}{return} is not \texorpdfstring{\texttt{print}}{print}}

\code{print} shows a \omterm{def:g6:variables-and-input:value}{value} to the person in front of the screen. \code{return}
gives a \omterm{def:g6:variables-and-input:value}{value} to the \omterm{def:g4:programs-and-operating-system:program}{program}. A \omterm{def:g7:functions:function}{function} that prints its result instead of
returning it shows the right number, but the \omterm{def:g4:programs-and-operating-system:program}{program} cannot use it:

\Needspace*{9\baselineskip}
\omcode{code/grade-7/01-functions/listings/print_not_return.py}{1}{6}{Printing
instead of returning.}\label{lst:g7:functions:print-not-return}

\Needspace*{11\baselineskip}
\omoutput{code/grade-7/01-functions/outputs/print_not_return_run.txt}{12 is
shown, then the addition fails.}\label{out:g7:functions:print-not-return}

A \omterm{def:g7:functions:function}{function} without \code{return} gives back the special \omterm{def:g6:variables-and-input:value}{value} \code{None},
which means ``nothing''. Adding 1 to nothing makes no sense: the \omterm{def:g6:first-python-program:error}{error message}
names the type of \code{None}, \code{NoneType}. The rule of thumb: a \omterm{def:g7:functions:function}{function}
that \emph{computes} something returns it, and the caller decides whether to
print it.

\section{Local variables}

\begin{definition}[Local variable]\label{def:g7:functions:local-variable}
A \emph{local variable}\index{local variable} is a \omterm{def:g6:variables-and-input:value}{variable} created inside a
\omterm{def:g7:functions:function}{function}: a \omterm{def:g7:functions:parameter}{parameter}, or a name assigned in the body. It exists only while
the call runs, and the rest of the \omterm{def:g4:programs-and-operating-system:program}{program} cannot see it.
\end{definition}

\Needspace*{10\baselineskip}
\omcode{code/grade-7/01-functions/listings/local.py}{1}{7}{A \omterm{def:g7:functions:local-variable}{local variable} used
outside its \omterm{def:g7:functions:function}{function}.}\label{lst:g7:functions:local}

\Needspace*{11\baselineskip}
\omoutput{code/grade-7/01-functions/outputs/local_run.txt}{\code{result} was
local to \code{area}: after the call, it is gone.}\label{out:g7:functions:local}

\begin{omfigure}
\begin{tikzpicture}[x=1cm, y=1cm, font=\footnotesize,
    vb/.style={draw=black!60, fill=ocYellow!25, minimum height=0.5cm, inner sep=2pt, font=\scriptsize}]
  \foreach \x/\t/\on in {0/{before the call}/0, 4.6/{during \code{area(3, 4)}}/1, 9.2/{after the call}/0} {
    \draw[black!40, rounded corners=4pt] (\x,0) rectangle ++(4.2,2.4);
    \node[anchor=north] at (\x+2.1,-0.05) {\t};
    \ifnum\on=1
      \draw[ocRed!70, dashed, rounded corners=3pt] (\x+0.15,0.15) rectangle ++(3.9,1.75);
      \node[ocRed!80!black, anchor=north west, font=\scriptsize] at (\x+0.2,1.85) {inside \code{area}};
      \node[vb] at (\x+1.05,1.05) {\code{width} 3};
      \node[vb] at (\x+3.05,1.05) {\code{length} 4};
      \node[vb] at (\x+2.1,0.45) {\code{result} 12};
    \else
      \node[black!50] at (\x+2.1,1.0) {no local boxes};
    \fi
  }
\end{tikzpicture}

{\small Local boxes appear when the call starts and disappear when it
returns. Only the \omterm{def:g7:functions:return-value}{return value}, 12, survives.}
\end{omfigure}

\omterm{def:g7:functions:local-variable}{Local variables} are a good thing: a \omterm{def:g7:functions:function}{function} can use any names it likes
without spoiling the \omterm{def:g6:variables-and-input:value}{variables} of the rest of the \omterm{def:g4:programs-and-operating-system:program}{program}, and two \omterm{def:g7:functions:function}{functions}
can both have a \omterm{def:g6:variables-and-input:value}{variable} called \code{total} without mixing them up.

\section{Cutting a problem into functions}

A long task is easier to write, read and correct when it is cut into small
\omterm{def:g7:functions:function}{functions}, each with one job. This receipt printer has three: one computes the
price of a line, one prints a line, and one prints the whole receipt.

\omcode{code/grade-7/01-functions/listings/receipt.py}{1}{19}{A receipt in three
\omterm{def:g7:functions:function}{functions}.}\label{lst:g7:functions:receipt}

\Needspace*{8\baselineskip}
\omoutput{code/grade-7/01-functions/outputs/receipt.txt}{The receipt; the columns
line up.}\label{out:g7:functions:receipt}

In the f-strings, the sign \code{:} after a \omterm{def:g6:variables-and-input:value}{value} gives its layout:
\code{:10} gives a text 10 places, and \code{:5.2f} writes a number 5 places
wide with 2 decimals.

\code{show_line} prints, which seems to break the rule of thumb above. Printing
a line is its job, and it also returns the line's total, so that
\code{receipt} can add the totals up.

\begin{omfigure}
\begin{tikzpicture}[x=1cm, y=1cm, font=\footnotesize,
    fn/.style={draw=black!65, rounded corners=3pt, fill=#1, minimum width=2.6cm, minimum height=0.7cm}]
  \node[fn=black!6] (m) at (0,0) {main program};
  \node[fn=ocYellow!30] (r) at (0,-1.2) {\code{receipt()}};
  \node[fn=ocBlue!12] (s) at (0,-2.4) {\code{show_line(...)}};
  \node[fn=ocMint!25] (l) at (0,-3.6) {\code{line_total(...)}};
  \draw[-{Stealth[length=2mm]}, thick] (m) -- (r);
  \draw[-{Stealth[length=2mm]}, thick] (r) -- node[right] {3 calls} (s);
  \draw[-{Stealth[length=2mm]}, thick] (s) -- node[right] {1 call each} (l);
\end{tikzpicture}

{\small Who calls whom in \cref{lst:g7:functions:receipt}.}
\end{omfigure}

Each \omterm{def:g7:functions:function}{function} should be checked before it is used by the others. The method:
work out a few answers by hand, then call the \omterm{def:g7:functions:function}{function} with the same \omterm{def:g7:functions:parameter}{arguments}
and compare. Choose examples that reach every branch, and the special cases
(zero, two equal numbers).

\omcode{code/grade-7/01-functions/listings/check_examples.py}{1}{16}{Checking
two \omterm{def:g7:functions:function}{functions} on examples.}\label{lst:g7:functions:check}

\Needspace*{9\baselineskip}
\omoutput{code/grade-7/01-functions/outputs/check_examples.txt}{Every answer
agrees.}\label{out:g7:functions:check}

\begin{remark}[Programs that are also toolboxes]\label{rem:g7:functions:main}
A \omterm{def:g3:files-and-folders:file}{file} of \omterm{def:g7:functions:function}{functions} can be used by another \omterm{def:g4:programs-and-operating-system:program}{program}, which takes its \omterm{def:g7:functions:function}{functions}
with \code{from area import area}, as we took \code{forward} from
\code{turtle}. The lines of the \omterm{def:g3:files-and-folders:file}{file} that are not in a \omterm{def:g7:functions:function}{function} then run as
well, which is not wanted. Python programmers therefore put them under
\code{if __name__ == "__main__":}. Those lines run when the \omterm{def:g3:files-and-folders:file}{file} itself is run,
and not when another \omterm{def:g4:programs-and-operating-system:program}{program} takes its \omterm{def:g7:functions:function}{functions}:

\omcode{code/grade-7/01-functions/listings/guarded.py}{1}{6}{The main part
of a \omterm{def:g3:files-and-folders:file}{file}, guarded.}\label{lst:g7:functions:guarded}

\omoutput{code/grade-7/01-functions/outputs/guarded.txt}{Run directly, the \omterm{def:g3:files-and-folders:file}{file}
prints its example.}\label{out:g7:functions:guarded}
\end{remark}

\section{Exercises}

\begin{exercise}[$\star$]\label{exo:g7:functions:1}
What does this \omterm{def:g4:programs-and-operating-system:program}{program} print?

\omcode{code/grade-7/01-functions/exercises/two_calls.py}{1}{7}{Two
calls.}\label{lst:g7:functions:two-calls}
\end{exercise}

\begin{exercise}[$\star$]\label{exo:g7:functions:2}
With \code{def f(x): return 2 * x + 1}, what does \code{print(f(2) + f(3))}
print?
\end{exercise}

\begin{exercise}[$\star$]\label{exo:g7:functions:3}
What appears on the screen?

\omcode{code/grade-7/01-functions/exercises/print_and_return.py}{1}{7}{Printing
and returning.}\label{lst:g7:functions:print-and-return}
\end{exercise}

\begin{exercise}[$\star$]\label{exo:g7:functions:4}
Write a \omterm{def:g7:functions:function}{function} \code{perimeter(width, height)} that returns the perimeter of a
rectangle.
\end{exercise}

\begin{exercise}[$\star\star$]\label{exo:g7:functions:5}
Tom's \omterm{def:g7:functions:function}{function} \code{average(a, b)} ends with \code{print((a + b) / 2)}. His
\omterm{def:g4:programs-and-operating-system:program}{program} then does \code{print(average(8, 12) * 2)} and stops with a
\code{TypeError}. Explain, and fix the \omterm{def:g7:functions:function}{function}.
\end{exercise}

\begin{exercise}[$\star\star$]\label{exo:g7:functions:6}
With the \omterm{def:g7:functions:function}{function} \code{area} of \cref{lst:g7:functions:area}, Noor calls
\code{area(4, 3)} for a room 3 wide and 4 long. Is the answer wrong? Now take
\code{introduce} of \cref{lst:g7:functions:two-params}: why does the order
matter there and not for \code{area}?
\end{exercise}

\begin{exercise}[$\star\star$]\label{exo:g7:functions:7}
Write a \omterm{def:g7:functions:function}{function} \code{is_even(n)} that returns \code{True} or \code{False}.
\end{exercise}

\begin{exercise}[$\star\star$]\label{exo:g7:functions:8}
Write a \omterm{def:g7:functions:function}{function} \code{minutes_to_hours(m)} that returns a number of minutes
in hours (90 gives 1.5) and prints nothing.
\end{exercise}

\begin{exercise}[$\star\star$]\label{exo:g7:functions:9}
How many times does the body of \code{square} run? What does the \omterm{def:g4:programs-and-operating-system:program}{program} print?

\omcode{code/grade-7/01-functions/exercises/loop_calls.py}{1}{9}{A call inside a
\omterm{def:g4:algorithms-repeat-and-decide:loop}{loop}.}\label{lst:g7:functions:loop-calls}
\end{exercise}

\begin{exercise}[$\star\star\star$]\label{exo:g7:functions:10}
Write a \omterm{def:g7:functions:function}{function} \code{price(age)} that returns the cinema ticket price of
\cref{lst:g6:if-else:ticket}, and check it on the four ages of
\cref{lst:g6:if-else:boundaries}.
\end{exercise}

\begin{exercise}[$\star\star\star$]\label{exo:g7:functions:11}
To check \code{largest3} of \cref{lst:g7:functions:largest}, which calls would
you make so that the largest number is once in each of the three places? Add
the cases with equal numbers. How many calls is that, at least?
\end{exercise}

\begin{exercise}[$\star\star\star$]\label{exo:g7:functions:12}
Write a \omterm{def:g7:functions:function}{function} \code{digit_sum(n)} that returns the sum of the digits of
\code{n}, and check it on five examples worked out by hand first.
\end{exercise}

\section{Problem: The Check Digit of a Barcode}

\begin{problem}[{Weekend problem --- catching typing mistakes}]\label{pb:g7:functions:1}
The 13-digit number under a barcode ends with a \emph{check digit}. It is
chosen so that the \emph{weighted sum} of the 13 digits is a multiple of 10:
counting from the right, the digits are multiplied by 1, 3, 1, 3, \dots\ in
turn, and the products are added. A cashier who types a number with a mistake
usually gets a weighted sum that is not a multiple of 10, and the till refuses
it.
% ledger: g7:ean13-rule

\needspace{4\baselineskip}\medskip\noindent\textbf{Part I --- By hand.} The number \code{5901234123457}
is a valid example.
\begin{enumerate}
  \item Compute its weighted sum. Is it a multiple of 10?
  \item The first twelve digits of another number are \code{400638133393}.
    Find its check digit.
  \item The cashier types \code{5901234123487}: one digit is wrong. Show that
    the till refuses it.
  \item Why is the weight of the check digit 1, and not 3, convenient for
    computing it?
\end{enumerate}

\needspace{4\baselineskip}\medskip\noindent\textbf{Part II --- Functions.} The number is kept as a whole
number; its digits are taken off from the right with \texttt{\% 10} and
\code{// 10}. The text in triple quotes under a \code{def} line is a
\emph{docstring}: a sentence that says what the \omterm{def:g7:functions:function}{function} does.

\omcode{code/grade-7/01-functions/listings/barcode.py}{1}{24}{Three \omterm{def:g7:functions:function}{functions} on
barcode numbers.}\label{lst:g7:functions:barcode}

\begin{enumerate}[resume]
  \item In \code{weighted_sum}, explain line 8: why does \code{4 - weight}
    switch between 1 and 3?
  \item In \code{check_digit}, why multiply \code{first12} by 10 before
    computing the weighted sum? Why is there a second \texttt{\% 10} at the end?
  \item Check \code{check_digit} on your answer to question 2, and
    \code{is_valid} on the numbers of questions 1 and 3.
  \item Write down how you would check \code{check_digit} on examples, in the
    style of \cref{lst:g7:functions:check}.
\end{enumerate}

\needspace{4\baselineskip}\medskip\noindent\textbf{Part III --- Which mistakes are caught?}

\omcode{code/grade-7/01-functions/listings/one_digit.py}{1}{16}{Every one-digit
change of a valid number.}\label{lst:g7:functions:one-digit}

\begin{enumerate}[resume]
  \item How many numbers does \cref{lst:g7:functions:one-digit} try? Run it:
    how many are refused?
  \item Explain why changing one digit with weight 3 always changes the
    weighted sum by an amount that is not a multiple of 10.
  \item Another common mistake is to swap two neighbouring digits. One of them
    has weight 1 and the other weight 3. If they are $a$ and $b$, by how much
    does the weighted sum change?
  \item The \omterm{def:g4:programs-and-operating-system:program}{program} \code{swaps.py} tries all 90 pairs of different digits $a,
    b$. Run it: how many swaps are caught, and which are missed?

\omcode{code/grade-7/01-functions/listings/swaps.py}{1}{11}{Every swap of two
neighbouring digits.}\label{lst:g7:functions:swaps}

  \item State the result.
\end{enumerate}
\end{problem}
